%% =================================================================
%% Chen et al. Measurement 264 (2026) 120300.
%% Reconstruction of page 3 (printed p. 3).
%% Contents: Fig. 4 (closed-loop feedback flow chart + Z0...Zn probe
%% descent sequence), Fig. 5 (a-d: AB glue hardness samples, Shore
%% hardness tester, compression-depth map, compression-depth vs
%% Shore-hardness fit). One embedded equation in Fig. 5 caption
%% ($Y = 6.23 - 1.62\,X^1 + 1.15\,X^2$, $R^2 = 0.99596$).
%% No numbered equations on this page.
%% =================================================================

\documentclass[11pt]{article}

\usepackage[margin=0.85in]{geometry}
\usepackage{amsmath, amssymb}
\usepackage{mathrsfs}
\usepackage{xcolor}
\usepackage{fancyhdr}

\pagestyle{fancy}
\fancyhf{}
\lhead{\small Z.~Chen et al.}
\rhead{\small Measurement 264 (2026) 120300}
\cfoot{\small 3 $\to$ \arabic{page}}
\renewcommand{\headrulewidth}{0.4pt}

\setcounter{page}{1}
\setlength{\parskip}{6pt}
\setlength{\parindent}{0pt}

\begin{document}

% ---------- Fig. 4 placeholder ----------
\begin{center}
\fbox{\begin{minipage}[c][6cm][c]{0.95\textwidth}
\centering
\textbf{Fig. 4. Flow chart of the closed-loop feedback control system.}\\[6pt]
\textcolor{gray}{\footnotesize
[main flow on the left: Planning scan path $\to$ PLC $\to$ ``The X
and Y axes move'' $\to$ ``Z-axis descent'' $\to$ decision diamond
``The pressure sensor reaches the first threshold'' (NO loops back
to Z-axis descent, YES feeds the grating-sensor position reader) $\to$
``The pressure sensor reaches the second threshold'' $\to$ \dots
$\to$ ``The pressure sensor reaches the final threshold'' $\to$
``Z-axis stops and resets.'' A left-hand feedback loop
($(Z_0, Z_1, \ldots, Z_n)$) carries position data from the grating
sensor back up into the PLC. On the right, a visual strip shows the
probe descending through layers of a stratified soft sample at
successive depths $Z_0$, $Z_1$, $\ldots$, $Z_n$.]}
\end{minipage}}
\end{center}

\vspace{0.8em}

% ---------- Fig. 5 placeholder ----------
\begin{center}
\fbox{\begin{minipage}[c][7cm][c]{0.95\textwidth}
\centering
\textbf{Fig. 5.} Softness recognition of the tactile tomography
system.\\[4pt]
\textcolor{gray}{\footnotesize
\textbf{(a)} 3 $\times$ 3 grid of nine AB glue samples, labeled with
Shore hardnesses 3, 7.5, 12.5, 17.5, 25, 30, 35, 50, and 60 HC.\\[3pt]
\textbf{(b)} Shore hardness tester (LX-C), photograph of the
bench-top analog dial instrument with inset views of LX-A, LX-C, and
LX-D indenter tips.\\[3pt]
\textbf{(c)} Compression depth distribution over the 3$\times$3 AB
glue grid as scanned by the tactile tomography system. X, Y in mm;
Z color scale 0 (dark blue) to 6 mm (dark red), showing compression
depth decreasing as Shore hardness increases.\\[3pt]
\textbf{(d)} Relation between compression depth and Shore hardness.
Measurement points (black squares) with a quadratic fit (red curve):
$Y = 6.23 - 1.62\,X^{1} + 1.15\,X^{2}$,\ $R^{2} = 0.99596$.
Compression depth axis: 0--7 mm; Shore hardness axis: 0--70 HC.]}
\end{minipage}}
\end{center}

\vspace{0.4em}

\textbf{Fig. 5.}\ Softness recognition of the tactile tomography
system. (a) AB glues with Shore hardnesses of 3, 7.5, 12.5, 17.5,
25, 30, 35, 50, and 60 HC. (b) Shore hardness tester (LX-C).
(c) Compression depth distribution when the tactile tomography
system scanned the AB glue samples with different Shore hardnesses
(scanned by the tactile tomography system). (d) Relation between
compression depth and Shore hardness.

\end{document}
